\documentclass[11pt]{article}
\usepackage[margin=1in]{geometry}
\usepackage{amsmath,amssymb,amsthm}
\usepackage[hidelinks]{hyperref}

\newtheorem{theorem}{Theorem}
\newtheorem{lemma}[theorem]{Lemma}
\newtheorem{corollary}[theorem]{Corollary}
\newtheorem{proposition}[theorem]{Proposition}
\theoremstyle{remark}
\newtheorem{remark}[theorem]{Remark}

\newcommand{\Tr}{\operatorname{Tr}}
\renewcommand{\Re}{\operatorname{Re}}
\newcommand{\norm}[1]{\lVert #1\rVert}

\title{Squared partial traces and the nuclear norm:\\ a multipartite inequality and the first-order Werner bound}
\author{Ansh Mishra \and Aryan Senthilkumar}
\date{October 2026}

\begin{document}
\maketitle

\begin{abstract}
We prove that every operator $Z$ on a finite-dimensional tensor product $H_1\otimes\cdots\otimes H_n$ satisfies
\[
\sum_{i=1}^n \norm{\Tr_i Z}_2^2\;\le\;(n-1)\norm{Z}_2^2+\norm{Z}_1^2 ,
\]
where $\Tr_i$ is the partial trace over the $i$th factor. The inequality is sharp. It needs no positivity and no rank
assumption. For $n=2$ it follows from a recent bound of Costa Rico and Wolf; for $n\ge3$ we believe it is new.

For operators of rank at most two it gives $\sum_i\norm{\Tr_iC}_2^2\le(n+1)\norm{C}_2^2$. This is Assumption~(i) of
Chattopadhyay and Sarkar from 2006, i.e.\ the first-order part of the $n$-copy distillability problem for Werner
states at the endpoint $\alpha=1/2$. Iterating the inequality gives all of their assumed bounds for partial traces over
$m$ sites. For density matrices it is a Han-type inequality for the linear (Tsallis-2) entropy,
$\sum_i S_L(\rho_{\hat\imath})\ge(n-1)S_L(\rho)$, which holds even though strong subadditivity fails for this entropy.

An induction on $n$ reduces the inequality to a single statement about block operators. We prove that statement with
a balanced factorisation and second compound matrices. The method does not reach the endpoint itself: every rank-free
relaxation of the endpoint inequality we tried fails already for two or three copies. In the remarks we split the
multiplicative (``no collective gain'') form of the problem into $n$ inequalities, of which ours is the first and the
endpoint the last. For three copies the other two follow from the two-copy theorem and Tomar's computer-assisted
three-copy theorem. So, given that theorem, the multiplicative form holds for three copies in every dimension; for local dimension at most four this was proved earlier by MacCormack and Bao, through an identity of the same kind. We also prove that the
two-copy inequality has a margin linear in the thickness of the co-range, which answers a question of Tomar, and that
for normal operators it holds with the sharp constant $1-1/\sqrt2$, better than the $\frac14$ Tomar conjectured.
\end{abstract}

\section{Introduction}

Let $F$ be the swap on $\mathbb C^d\otimes\mathbb C^d$. For $\alpha\in[-1,1]$ let $\rho_\alpha\propto I-\alpha F$ be the
Werner state. It has a non-positive partial transpose (NPT) for $\alpha>1/d$ and is one-copy distillable for
$\alpha>1/2$. Whether the states with $1/d<\alpha\le1/2$ can be distilled from some number of copies is the NPT
bound-entanglement problem \cite{DSSTT,DCLB,HRZ}.

The two-copy case was settled in 2026 by several groups \cite{FGP,SongChen,FHPV,BGH}. Three copies have since been
settled at the endpoint with computer assistance: for $d=3$ and $d=4$ by exact certificates \cite{AnonD3,LiuD4}, which MacCormack and Bao replayed, formalising $d=3$ in Lean and proving the sharp three-copy curve for $d\le4$ \cite{MB}; and in every dimension by Tomar \cite{Tomar}. Earlier, Wu and Zou proved the three-copy
endpoint inequality for positive semidefinite and for normal operators of rank two, and for operators whose range and
co-range lie in one plane \cite{WuZou}. The general $n$-copy
question remains open.

By \cite[Thm.~1]{CR23} (see also \cite[Thm.~V.1]{QGWG}; restated as \cite[Thm.~1]{SongChen}), $\rho_{1/2}$ is $n$-copy undistillable if and only if
\begin{equation}\label{eq:qn}
q_n(C)\;=\;\sum_{J\subseteq\{1,\dots,n\}}\Big(-\tfrac12\Big)^{|J|}\,\norm{\Tr_J C}_2^2\;\ge\;0
\end{equation}
for every operator $C$ on $(\mathbb C^d)^{\otimes n}$ of rank at most two. Here $\Tr_J$ is the partial trace over the
factors in $J$.

Long before these results, Chattopadhyay and Sarkar \cite{CS} proposed bounds for Schmidt-rank-two vectors $\psi$ in
$(\mathbb C^d\otimes\mathbb C^d)^{\otimes n}$. Let $\psi$ be the vectorisation of $C$, normalised so that
$\norm{C}_2=1$. With the unnormalised projector $P^+=\sum_{j,k}|jj\rangle\langle kk|$ (trace $d$) one has
$\langle\psi|P^+_J\otimes I|\psi\rangle=\norm{\Tr_JC}_2^2$, where $P^+_J$ acts on the pairs in $J$. Their bounds then read
\begin{equation}\label{eq:CS}
S_m(C):=\sum_{|J|=m}\norm{\Tr_JC}_2^2\;\le\;\big[\tbinom{n-1}{m}+2\tbinom{n-1}{m-1}\big]\norm{C}_2^2,\qquad m=1,\dots,n,
\end{equation}
for every $C$ of rank at most two.

Chattopadhyay and Sarkar presented these bounds as assumptions supported by numerics. Earlier versions of their paper
claimed more; Pankowski, Piani and the Horodeckis later noted that the proof had gaps \cite{PPHH}. Two cases are now
known.
\begin{itemize}
\item The case $n=2$ is $\norm{\Tr_1C}_2^2+\norm{\Tr_2C}_2^2\le3\norm{C}_2^2$ \cite[Eq.~(32)]{CRW}. It also follows from
the sharper rank-two inequalities of \cite{FHPV,BGH}.
\item The case $m=n-1$ involves only one-site marginals. It follows from Corollary~2 of Costa Rico and Wolf \cite{CRW}:
$\sum_i\norm{\Tr_{\hat\imath}M}_2^2\le(1+r(n-1))\norm{M}_2^2$ for operators $M$ of rank $r$.
\end{itemize}
Their separate two-copy assumption (Eq.~(3) of \cite{CS}) follows from the rank-two bipartite bounds of
\cite{FHPV,BGH}.

The case $m=1$, $\sum_i\norm{\Tr_iC}_2^2\le(n+1)\norm{C}_2^2$, is the first-order part of \eqref{eq:qn}: it is the
derivative at $t=0$ of the conjectured identity
$\min_C\sum_J(-t)^{|J|}\norm{\Tr_JC}_2^2=(1-2t)(1-t)^{n-1}\norm{C}_2^2$ for rank-two $C$.

Our main result proves all of \eqref{eq:CS}, in particular the first-order case $m=1$ for every $n\ge3$, as a
consequence of one inequality that holds for every operator.

Section~\ref{sec:remarks} contains further results; we mention three. The conjectured identity above splits into $n$ inequalities,
of which ours is the first and \eqref{eq:qn} the last. For three copies the middle one follows from the two-copy
theorem and Tomar's three-copy theorem. Finally, for two copies we prove that $q_2(C)$ is bounded below by a constant
times the ``thickness'' of the co-range of $C$ (Proposition~\ref{prop:thick}), which answers a question of
Tomar~\cite{Tomar2}; when the range and co-range of $C$ lie in one plane, for instance for normal $C$, we obtain the
sharp constant $1-1/\sqrt2$, which beats the $\frac14$ he conjectured. Tomar proved this constant on a different class
of planes \cite[Thm.~7.2]{Tomar2}.

Partial-trace norm inequalities have a long history.
\begin{itemize}
\item Audenaert proved subadditivity of $q$-entropies for $q>1$ \cite{Aud}.
\item Rastegin studied unitarily invariant norms of partial traces \cite{Ras}.
\item Costa Rico and Wolf extended Audenaert's inequality to general matrices and several tensor factors \cite{CRW}.
Their bounds, such as $\sum_i\norm{M_i}_p\le(n-1)\norm{M}_1+\norm{M}_p$ (their Prop.~9), concern the one-site marginals
$M_i$. In the notation above they control $S_{n-1}$, whereas the first-order term is $S_1$. For $n=2$ the two kinds of
marginal coincide. There Prop.~9 with $p=2$, followed by an elementary squaring step, gives the case $n=2$ of
Theorem~\ref{thm:main} below.
\item The two-copy proofs rely on sharp bipartite inequalities for operators of bounded rank \cite{FHPV,BGH}.
\item Wu and Zou proved a sharp rank-free bipartite inequality,
$\norm{\Tr_KX}_2^2+\norm{\Tr_LX}_2^2\le\norm{X}_1^2+|\Tr X|^2$ \cite[Thm.~1.1]{WuZou}. Neither this nor our $n=2$ case
implies the other: the identity favours ours, and traceless operators favour theirs.
\item The first step of our block lemma is the second-compound estimate of Wu and Zou \cite[Lemma~5.1]{WuZou}, in a
homogeneous form.
\end{itemize}
For $n\ge3$ we are not aware of a previous statement of Theorem~\ref{thm:main}, of Lemma~\ref{lem:R}, or of the bounds
\eqref{eq:CS} with $m<n-1$.

\section{Results}

Throughout, $H=H_1\otimes\cdots\otimes H_n$ with $\dim H_i<\infty$. We write $\norm{\cdot}_2$ for the
Hilbert--Schmidt norm, $\norm{\cdot}_1$ for the trace (nuclear) norm and $\langle A,B\rangle=\Tr(A^*B)$.

\begin{theorem}\label{thm:main}
For every operator $Z$ on $H$,
\begin{equation}\label{eq:main}
\sum_{i=1}^n\norm{\Tr_iZ}_2^2\;\le\;(n-1)\norm{Z}_2^2+\norm{Z}_1^2 .
\end{equation}
Equality holds for $Z=e^{i\theta}P\otimes|c\rangle\langle c|$ with $P\ge0$ on one factor and $c$ a product unit vector
on the others.
\end{theorem}

Since $\norm{C}_1^2\le\operatorname{rank}(C)\,\norm{C}_2^2$, Theorem~\ref{thm:main} has an immediate consequence.

\begin{corollary}\label{cor:rank}
If $\operatorname{rank}C\le r$, then $\sum_i\norm{\Tr_iC}_2^2\le(n+r-1)\norm{C}_2^2$. For $r=2$ this is \eqref{eq:CS}
with $m=1$, for every $n$ and every local dimension, including nonnormal $C$.
\end{corollary}

The bound is sharp for every $r$: take $C=I_U\otimes|c\rangle\langle c|$ with $U$ an $r$-dimensional subspace of one
factor and $c$ a product vector. It is the first-order term of the natural rank-$r$ form of the multiplicativity
conjecture,
$\min\{\sum_J(-t)^{|J|}\norm{\Tr_JC}_2^2:\operatorname{rank}C\le r,\ \norm{C}_2=1\}=(1-rt)(1-t)^{n-1}$, whose
derivative at $t=0$ is $-(n+r-1)$. The case $r=2$, $t=1/2$ is \eqref{eq:qn}, i.e.\ the statement that the map
$X\mapsto\Tr(X)I-X/2$ is tensor-stably $2$-positive.

In the language of Schmidt-number witnesses \cite{CRW}, Corollary~\ref{cor:rank} says the following. On
$(\mathbb C^d\otimes\mathbb C^d)^{\otimes n}$ the operator
\[
W_r=(n+r-1)\,I-\sum_{i=1}^nP^+_{A_iB_i}\otimes I_{\rm rest}
\]
is nonnegative on every vector of Schmidt rank at most $r$ across $A|B$. For $r<d$ it is negative on a vector of
Schmidt rank $r+1$. Equivalently, $X\mapsto(n+r-1)\Tr(X)\,I-\sum_i\Tr_{\hat\imath}(X)\otimes I_{\hat\imath}$ is
$r$-positive but not $(r+1)$-positive. For $r=2$ this is the operator of \cite{CS}. Costa Rico and Wolf give the
analogous family in which only one copy is left untraced \cite{CRW}.

\begin{corollary}\label{cor:allm}
For every operator $Z$ on $H$ and every $0\le m\le n$,
\[
S_m(Z)\le\tbinom{n-1}{m}\norm{Z}_2^2+\tbinom{n-1}{m-1}\norm{Z}_1^2,
\]
with $\binom{n-1}{-1}=0$. Equality holds for $Z=P\otimes|c\rangle\langle c|$ as in Theorem~\ref{thm:main}. In particular
all the bounds \eqref{eq:CS} hold.
\end{corollary}

\begin{corollary}\label{cor:han}
For every density matrix $\rho$ on $H$, with $\rho_{\hat\imath}=\Tr_i\rho$ and $S_L(\sigma)=1-\Tr\sigma^2$,
\[
\sum_{i=1}^n\Tr\rho_{\hat\imath}^2\le(n-1)\Tr\rho^2+1,
\qquad\text{equivalently}\qquad
\sum_{i=1}^nS_L(\rho_{\hat\imath})\ge(n-1)S_L(\rho).
\]
\end{corollary}

For the von Neumann entropy this is Han's inequality. Its usual proof uses strong subadditivity, which fails for $S_L$.
The case $n=2$ is subadditivity of $S_L$ \cite{Aud}. For diagonal $\rho$, Corollary~\ref{cor:han} is the case $q=2$ of
the classical Han inequality for Tsallis entropies \cite{FYK}.

\begin{corollary}\label{cor:map}
Let $A_0$ be an auxiliary system. For states $\sigma$ on $A_0\otimes H$ write $\sigma_0$ and $\sigma_{0i}$ for the
marginals on $A_0$ and on $A_0\otimes H_i$. The map
\[
\Lambda(\sigma)=\Tr(\sigma)\,I+(n-1)\,\sigma_0\otimes I_H-\sum_{i=1}^n\sigma_{0i}\otimes I_{H_{\hat\imath}}
\]
is positive. Consequently every state on $A_0\otimes H$ satisfies $\sum_i\Tr\sigma_{0i}^2\le(n-1)\Tr\sigma_0^2+1$.
\end{corollary}

For $n=1$, $\Lambda$ is the reduction map. For $n\ge2$ one checks numerically (not scripted here) that $\Lambda$ is not decomposable, so its
positivity is not a consequence of positivity and partial transposition alone.

\section{The block lemma}

\begin{lemma}\label{lem:R}
Let $K$ be finite-dimensional and let $Z=\sum_{a,b=1}^d|a\rangle\langle b|\otimes Z_{ab}$ be an operator on
$\mathbb C^d\otimes K$ with blocks $Z_{ab}$ on $K$. Then
\begin{equation}\label{eq:R}
\sum_{a,b}\norm{Z_{ab}}_1^2+\Big\|\sum_aZ_{aa}\Big\|_2^2\;\le\;\sum_{a,b}\norm{Z_{ab}}_2^2+\norm{Z}_1^2 .
\end{equation}
\end{lemma}

Equality holds for $Z=|v\rangle\langle v|\otimes W$ with any unit $v$ and any $W$.

\begin{proof}
If $Z=0$ there is nothing to prove, so let $r=\operatorname{rank}Z\ge1$. Let $Z=U\Sigma V^*$ be a compact singular value
decomposition, with $\Sigma>0$ diagonal of size $r$. Put $P=U\Sigma^{1/2}$ and $Q=V\Sigma^{1/2}$. Then $Z=PQ^*$ and
$P^*P=Q^*Q=\Sigma$.

Let $P_a=(\langle a|\otimes I_K)P$ and $Q_b=(\langle b|\otimes I_K)Q$. These are maps from $\mathbb C^r$ to $K$, and
$Z_{ab}=P_aQ_b^*$. Put $R_a=P_a^*P_a$ and $S_b=Q_b^*Q_b$. Then $\sum_aR_a=\sum_bS_b=\Sigma$. Write
\[
\alpha_a=\Tr R_a,\quad p_a=\Tr R_a^2,\quad \beta_b=\Tr S_b,\quad q_b=\Tr S_b^2,
\]
so that $0\le p_a\le\alpha_a^2$ and $0\le q_b\le\beta_b^2$. Then $\norm{Z}_1=\Tr\Sigma=\sum_a\alpha_a=\sum_b\beta_b$, and
therefore $\norm{Z}_1^2=\sum_{a,b}\alpha_a\beta_b$.

\emph{Step 1 (off-diagonal and diagonal blocks alike).}
\begin{itemize}
\item For any matrix $M$ with singular values $s_j$,
$\norm{M}_1^2-\norm{M}_2^2=2\sum_{j<k}s_js_k=2\norm{\wedge^2M}_1$.
Here $\wedge^2$ is the second compound: for a linear map $A$, $\wedge^2A$ is the induced map on antisymmetric tensors.
It vanishes when the domain or the range has dimension one.
\item The compound is multiplicative and commutes with adjoints, also for rectangular maps (the Cauchy--Binet formula).
So $\wedge^2Z_{ab}=(\wedge^2P_a)(\wedge^2Q_b)^*$.
\item By H\"older's inequality, $\norm{XY^*}_1\le\norm{X}_2\norm{Y}_2$.
\item Moreover $\norm{\wedge^2P_a}_2^2=\Tr\wedge^2(P_a^*P_a)=\tfrac12(\alpha_a^2-p_a)$, and likewise for $Q_b$.
\end{itemize}
Hence
\begin{equation}\label{eq:s1}
\norm{Z_{ab}}_1^2-\norm{Z_{ab}}_2^2\le\sqrt{(\alpha_a^2-p_a)(\beta_b^2-q_b)}\le\alpha_a\beta_b-\sqrt{p_aq_b}.
\end{equation}
The first inequality is \cite[Lemma~5.1]{WuZou} in homogeneous form. For the second, both sides are nonnegative, and after
squaring it reduces to $\alpha_a^2q_b+\beta_b^2p_a\ge2\alpha_a\beta_b\sqrt{p_aq_b}$.

\emph{Step 2 (the diagonal blocks).} By Cauchy--Schwarz,
\[
\norm{Z_{aa}}_2^2=\Tr(Q_aP_a^*P_aQ_a^*)=\Tr(R_aS_a)\le\norm{R_a}_2\norm{S_a}_2=\sqrt{p_aq_a}.
\]
By the triangle inequality and Cauchy--Schwarz again,
\begin{equation}\label{eq:s2}
\Big\|\sum_aZ_{aa}\Big\|_2\le\sum_a(p_aq_a)^{1/4}\le\Big(\sum_a\sqrt{p_a}\Big)^{1/2}\Big(\sum_a\sqrt{q_a}\Big)^{1/2}.
\end{equation}

\emph{Step 3 (conclusion).} Let $\Delta$ be the right side of \eqref{eq:R} minus the left side. Using
$\norm{Z}_1^2=\sum_{a,b}\alpha_a\beta_b$, then \eqref{eq:s1} and \eqref{eq:s2},
\begin{align*}
\Delta&=\sum_{a,b}\big[\alpha_a\beta_b-\big(\norm{Z_{ab}}_1^2-\norm{Z_{ab}}_2^2\big)\big]-\Big\|\sum_aZ_{aa}\Big\|_2^2\\
&\ge\sum_{a,b}\sqrt{p_aq_b}-\Big(\sum_a\sqrt{p_a}\Big)\Big(\sum_a\sqrt{q_a}\Big)=0. \qedhere
\end{align*}
\end{proof}

\section{Proofs}

\begin{proof}[Proof of Theorem~\ref{thm:main}]
We use induction on $n$. For $n=1$ the claim is $|\Tr Z|^2\le\norm{Z}_1^2$.

For $n\ge2$ write $Z=\sum_{a,b}|a\rangle\langle b|\otimes Z_{ab}$ with respect to an orthonormal basis of $H_n$, with
blocks $Z_{ab}$ on $H_1\otimes\cdots\otimes H_{n-1}$. Then:
\begin{itemize}
\item $\Tr_iZ=\sum_{a,b}|a\rangle\langle b|\otimes\Tr_iZ_{ab}$ for $i<n$;
\item $\Tr_nZ=\sum_aZ_{aa}$;
\item $\norm{Z}_2^2=\sum_{a,b}\norm{Z_{ab}}_2^2$.
\end{itemize}
The induction hypothesis, applied to each block, gives
\[
\sum_{i<n}\norm{\Tr_iZ}_2^2=\sum_{a,b}\sum_{i<n}\norm{\Tr_iZ_{ab}}_2^2\le(n-2)\norm{Z}_2^2+\sum_{a,b}\norm{Z_{ab}}_1^2.
\]
Adding $\norm{\Tr_nZ}_2^2$ and applying Lemma~\ref{lem:R}, with $K=H_1\otimes\cdots\otimes H_{n-1}$ and $d=\dim H_n$,
gives \eqref{eq:main}. Local dimensions may differ, since Lemma~\ref{lem:R} allows any $d$ and $K$.

For the equality case, let $Z=P\otimes|c\rangle\langle c|$ with $P\ge0$ on the first factor. Then:
\begin{itemize}
\item $\norm{\Tr_1Z}_2^2=(\Tr P)^2=\norm{P}_1^2$;
\item $\norm{\Tr_iZ}_2^2=\norm{P}_2^2$ for $i\ge2$;
\item $\norm{Z}_2=\norm{P}_2$ and $\norm{Z}_1=\norm{P}_1$.
\end{itemize}
So both sides of \eqref{eq:main} equal $(n-1)\norm{P}_2^2+\norm{P}_1^2$.
\end{proof}

\begin{proof}[Proof of Corollary~\ref{cor:allm}]
We use induction on $m$, for all $n$ at once. For $m=0$ the claim is trivial, and for $m=1$ it is
Theorem~\ref{thm:main}. Let $m\ge2$, so $n\ge2$.

Every $J$ with $|J|=m$ satisfies
\[
\sum_{i\in J}\norm{\Tr_{J\setminus i}(\Tr_iZ)}_2^2=m\norm{\Tr_JZ}_2^2 .
\]
Hence $mS_m(Z)=\sum_{i=1}^nS_{m-1}(\Tr_iZ)$, where each $\Tr_iZ$ is an operator on $n-1$ factors.

Partial traces do not increase the trace norm. Combining this with the induction hypothesis on $n-1$ factors and
Theorem~\ref{thm:main} gives
\[
mS_m(Z)\le\tbinom{n-2}{m-1}\big[(n-1)\norm{Z}_2^2+\norm{Z}_1^2\big]+n\tbinom{n-2}{m-2}\norm{Z}_1^2 .
\]
To finish, use $(n-1)\binom{n-2}{m-1}=m\binom{n-1}{m}$ and $\binom{n-2}{m-1}+n\binom{n-2}{m-2}=m\binom{n-1}{m-1}$. The
second identity follows from $\binom{n-2}{m-1}=\frac{n-m}{m-1}\binom{n-2}{m-2}$.

For the equality case: for $Z=P\otimes|c\rangle\langle c|$, one has
$S_m(Z)=\binom{n-1}{m-1}\norm{P}_1^2+\binom{n-1}{m}\norm{P}_2^2$. For rank two, $\norm{C}_1^2\le2\norm{C}_2^2$ turns the
bound into \eqref{eq:CS}.
\end{proof}

\begin{proof}[Proof of Corollaries~\ref{cor:han} and \ref{cor:map}]
For $\rho\ge0$ we have $\norm{\rho}_1=\Tr\rho$; this gives Corollary~\ref{cor:han}.

For Corollary~\ref{cor:map} it suffices, by linearity, to take $\sigma=|\psi\rangle\langle\psi|$. For any vector
$\varphi$ and any set $T$ of factors, the swap trick gives
$\langle\varphi|\sigma_T\otimes I|\varphi\rangle=\norm{\Tr_T(|\psi\rangle\langle\varphi|)}_2^2$, where $\sigma_T$ is the
marginal on $T$ and $\Tr_T$ traces out $T$. Hence, with $Z=\Tr_{A_0}(|\psi\rangle\langle\varphi|)$,
\[
\langle\varphi|\Lambda(\sigma)|\varphi\rangle=\norm{\psi}^2\norm{\varphi}^2+(n-1)\norm{Z}_2^2-\sum_i\norm{\Tr_iZ}_2^2 .
\]
This is nonnegative by Theorem~\ref{thm:main}, since $\norm{Z}_1\le\norm{\psi}\norm{\varphi}$.

The last statement is $\Tr[\sigma\Lambda(\sigma)]\ge0$.
\end{proof}

\section{Remarks and open problems}\label{sec:remarks}
The numerical statements in this section (``in our searches'', ``numerically'', ``we found'') come from exploratory local optimisation. Unless a script in \texttt{code/} is named, they are not part of the scripted checks.


\emph{Other multipartite forms.} The bipartite inequality of Wu and Zou, with $|\Tr X|^2$ in place of $\norm{X}_2^2$,
does not extend in the same way. The bound
$\sum_i\norm{\Tr_iZ}_2^2\le(n-2)\norm{Z}_2^2+\norm{Z}_1^2+|\Tr Z|^2$ fails for $n=3$: a numerical search finds rank-two operators
on three qubits that violate it by about $2.5\%$ of $\norm{Z}_1^2$.

\emph{A sharper rank-two bound.} For $n=2$ the sharp rank-two inequality of \cite{FHPV,BGH} reads
$\norm{\Tr_1C}_2^2+\norm{\Tr_2C}_2^2\le2\norm{C}_2^2+\tfrac12|\Tr C|^2$. Numerically, the multipartite version
\[
\sum_{i=1}^n\norm{\Tr_iC}_2^2\le n\norm{C}_2^2+\tfrac12|\Tr C|^2\qquad(\operatorname{rank}C\le2)
\]
holds in our searches, with equality attained, for $n\le6$ ($d=2$), $n=3,\dots,6,8,10,12$ ($d=3$), $n=4,6$ ($d=4$) and
$n=3,5$ ($d=5$). For positive $C$ it
follows from Corollary~\ref{cor:han}, since then
$(\Tr C)^2\le2\Tr C^2$. It then holds for every normal $C$ of rank two. Write $C=\lambda_1uu^*+\lambda_2vv^*$ with
orthonormal $u,v$, put $a=|\lambda_1|$, $b=|\lambda_2|$, $r=\Re(\bar\lambda_1\lambda_2)$ and $C_+=a\,uu^*+b\,vv^*$.
Then $S_1(C)=a^2S_1(uu^*)+b^2S_1(vv^*)+2rc$ with $c=\sum_i\langle\Tr_iuu^*,\Tr_ivv^*\rangle\ge0$. If $r\le0$, the
rank-one bound $S_1(uu^*)\le n$ gives $S_1(C)\le n\norm{C}_2^2$. If $r>0$, put $\theta=r/(ab)\le1$. Then
$S_1(C)=(1-\theta)\,(a^2S_1(uu^*)+b^2S_1(vv^*))+\theta\,S_1(C_+)\le n(a^2+b^2)+\theta(a+b)^2/2$, and this is at most
$n\norm{C}_2^2+\tfrac12|\Tr C|^2=n(a^2+b^2)+\tfrac12(a^2+b^2+2r)$.

The bound also holds whenever the ranges of $C$ and $C^*$ together span at most two dimensions, that is, for
$C=VMV^*$ with $V:\mathbb C^2\to H$ an isometry and $M$ any $2\times2$ matrix. Write $C=Z_1+iZ_2$ with $Z_1=(C+C^*)/2$ and $Z_2=(C-C^*)/2i$. Both are
Hermitian of rank at most two. For Hermitian $X,Y$ the numbers $\langle\Tr_JX,\Tr_JY\rangle$ are real, so
$\norm{\Tr_JC}_2^2=\norm{\Tr_JZ_1}_2^2+\norm{\Tr_JZ_2}_2^2$ for every $J$. Hence
$F(C)=S_1(C)-n\norm{C}_2^2-\tfrac12|\Tr C|^2$ satisfies $F(C)=F(Z_1)+F(Z_2)\le0$ by the Hermitian case. The same
splitting shows that for such $C$ the endpoint inequality \eqref{eq:qn} reduces to Hermitian operators of rank two.
When $\operatorname{ran}C+\operatorname{ran}C^*$ has dimension three or four we do not have a proof; the Hermitian part
of $C$ can then have rank four, and $F$ can be positive on it.

\emph{Sharp forms with a flatness term.} Let $\sigma_1\ge\sigma_2$ be the singular values of a rank-two $C$. For two
factors the sharp rank-two inequality can be improved to
\[
\norm{\Tr_1C}_2^2+\norm{\Tr_2C}_2^2\le\norm{C}_2^2+\tfrac12\norm{C}_1^2+\tfrac12|\Tr C|^2,
\qquad\text{i.e.}\qquad q_2(C)\ge\tfrac14(\sigma_1-\sigma_2)^2 .
\]
Since $\norm{C}_1^2=\norm{C}_2^2+2\sigma_1\sigma_2$, this is stronger than the bound of \cite{FHPV,BGH} by
$\tfrac12(\sigma_1-\sigma_2)^2$. It is an equality for every $C=\rho\otimes|c\rangle\langle c|$ with $\rho\ge0$ of rank
two, not only for flat $\rho$. It follows from the proof in \cite{FHPV}. There $C=x_1y_1^*+x_2y_2^*$ with
$\langle x_i,x_j\rangle=\langle y_i,y_j\rangle=g_{ij}$ and $\langle y_i,x_i\rangle=\Tr C/2$. To build it, take a
singular value decomposition $C=U\Sigma V^*$, put $X_0=U\Sigma^{1/2}$ and $Y_0=V\Sigma^{1/2}$, and replace them by
$X_0Q$ and $Y_0Q$, where the unitary $Q$ makes both diagonal entries of $Q^*Y_0^*X_0Q$ equal to $\Tr C/2$. This is
possible because the numerical range of a $2\times2$ matrix contains half its trace. Then $G=(g_{ij})=Q^*\Sigma Q$, so
$g_{11}+g_{22}=\norm{C}_1$ and $\norm{C}_2^2=g_{11}^2+g_{22}^2+2|g_{12}|^2$. For this decomposition we have the identity
\[
\norm{C}_2^2+\tfrac12\norm{C}_1^2+\tfrac12|\Tr C|^2-\sum_{i=1,2}\norm{\Tr_iC}_2^2
=\tfrac12(g_{11}-g_{22})^2+4\sum_{(a,b)}\norm{\Pi(a\otimes b)}^2+\sum_{i=1,2}\norm{\Tr_i(x_2x_1^*-y_2y_1^*)}_2^2 .
\]
Here $\Pi=\frac14(I-F_1)(I-F_2)$ on two copies of $H_1\otimes H_2$, $F_i$ swaps factor $i$ between the copies, and
$(a,b)$ runs over $(x_1,y_1)$, $(x_2,y_2)$, $(x_1,x_2)$, $(y_1,y_2)$. Every term on the right is nonnegative.

To prove the identity, write $\Phi(Z)=(\Tr_1Z,\Tr_2Z)$, so that $\sum_i\norm{\Tr_iC}_2^2=\norm{\Phi(C)}^2$. The swap trick
gives $\langle\Phi(uv^*),\Phi(st^*)\rangle=\langle u\otimes t|F_1+F_2|v\otimes s\rangle$ for all vectors $u,v,s,t$. Since
$F_1F_2$ exchanges the two copies, this has two consequences:
\[
\langle\Phi(uv^*),\Phi(st^*)\rangle=\langle\Phi(tv^*),\Phi(su^*)\rangle,\qquad
\norm{\Phi(xy^*)}^2=\norm{x}^2\norm{y}^2+|\langle x,y\rangle|^2-4\norm{\Pi(x\otimes y)}^2 ,
\]
the second because $4\Pi=I-F_1-F_2+F_1F_2$. Now expand $\norm{\Phi(C)}^2$. By the first relation the cross term is
$2\operatorname{Re}\langle\Phi(y_2y_1^*),\Phi(x_2x_1^*)\rangle=\norm{\Phi(x_2x_1^*)}^2+\norm{\Phi(y_2y_1^*)}^2
-\norm{\Phi(x_2x_1^*-y_2y_1^*)}^2$. Applying the second relation to the four rank-one terms gives
$\norm{\Phi(C)}^2=g_{11}^2+g_{22}^2+\frac12|\Tr C|^2+2g_{11}g_{22}+2|g_{12}|^2-4\sum_{(a,b)}\norm{\Pi(a\otimes b)}^2
-\norm{\Phi(x_2x_1^*-y_2y_1^*)}^2$, which is the identity. For
rank $r$ the same steps, starting from \cite[Eq.~(65)]{FHPV}, give $\norm{C}_2^2+(1-\frac1r)\norm{C}_1^2+\frac1r|\Tr C|^2$
on the right, which is again an equality for $\rho\otimes|c\rangle\langle c|$ with $\rho\ge0$ of rank $r$.

For $n$ factors the analogues would be
\[
q_n(C)\ge2^{-n}(\sigma_1-\sigma_2)^2,\qquad
S_1(C)\le(n-1)\norm{C}_2^2+\tfrac12\big(\norm{C}_1^2+|\Tr C|^2\big)\qquad(\operatorname{rank}C\le2).
\]
Both are equalities for $\rho\otimes(cc^*)^{\otimes(n-1)}$ with $\rho\ge0$ of rank two. In our searches the first holds
for $3\le n\le10$ at $d=3$, and the second for $n=3,4,6,8,10$ at $d=3$, with spot checks at $d=2,4,5,6$. For rank one the
first is \cite[Prop.~8]{SongChen}. The second adds $\tfrac12(\sigma_1-\sigma_2)^2$ to the slack of the sharp rank-two
bound above, and it holds for normal $C$: in the argument for normal $C$ above, Theorem~\ref{thm:main} applied to
$C_+$ gives $S_1(C)\le n\norm{C}_2^2+2\max(r,0)$, while the right-hand side equals $n\norm{C}_2^2+ab+r$ and $|r|\le ab$.
We could not prove either statement in general for $n\ge3$: the block induction behind Theorem~\ref{thm:main} fails for
the second already at $n=3$, whichever factor is split off. Tomar's Conjecture~7.1 \cite{Tomar2} strengthens
\eqref{eq:qn} in a different direction, with a margin given by the third eigenvalue of the one-copy reductions of the
co-range. For two copies we prove a margin of that kind below.

The second statement has a natural version for rank at most $r$:
\[
S_1(C)\le(n-1)\norm{C}_2^2+\big(1-\tfrac1r\big)\norm{C}_1^2+\tfrac1r|\Tr C|^2 .
\]
It is an equality for $\rho\otimes(cc^*)^{\otimes(n-1)}$ with $\rho\ge0$ of rank $r$, and as $r\to\infty$ it becomes
Theorem~\ref{thm:main}. For $n=2$ it is the bound proved above. For $r=1$ it holds for every $n$: it then reads
$S_1(xy^*)\le(n-1)\norm{x}^2\norm{y}^2+|\langle y,x\rangle|^2$, that is
$\sum_i\rho_i\otimes I\le(n-1)(\Tr\rho)\,I+\rho$ for $\rho=xx^*$, where $\rho_i$ is the marginal on factor $i$ and $I$
acts on the other factors. For $n=2$ this is \cite[Prop.~1]{CR23} (see also \cite[Prop.~11]{CRW} and \cite[Lemma~6.4]{FHPV}). For larger $n$, apply the statement for $n-1$
factors to the marginal of $\rho$ on the first $n-1$ factors, add $\rho_n\otimes I$, and use the case $n=2$ for the cut
between the first $n-1$ factors and the last. In our searches the rank-$r$ statement also holds for $r=3,4$ at
$n=3,4$.

\emph{A thickness margin for two copies.} For a two-dimensional subspace $V\subseteq H_1\otimes H_2$ and $l=1,2$, let
$\tau_l(V)$ be the sum of all but the two largest eigenvalues of the reduced operator of the projector $P_V$ on factor
$l$. It is at least the third eigenvalue $\lambda_3^{(l)}(V)$ used by Tomar, and it vanishes exactly when the local
support of $V$ at factor $l$ has dimension at most two.

\begin{proposition}\label{prop:thick}
Let $C$ have rank at most two on $H_1\otimes H_2$, and let $V$ be a two-dimensional subspace that contains the range or
the co-range of $C$. Then
\[
q_2(C)\ge\frac{\min\{\tau_1(V),\tau_2(V)\}}{48-8\sqrt2}\,\norm{C}_2^2 .
\]
\end{proposition}

In the notation of \cite{Tomar2} this gives $m_2(V)\ge\frac1{37}\min_l\lambda_3^{(l)}(V)$ for every plane $V$. That is
the case $k=2$ of Tomar's Conjecture~7.1 with a smaller constant: he conjectures $\frac14$, and in our searches the
best constant at $d=3$ is about $0.29$, in line with Tomar's numerics. For two copies no linear bound of this kind was known
\cite[Sec.~8]{Tomar2}. In our searches the margin is even linear in the larger tail $\max_l\tau_l(V)$, with infimum
$(3-\sqrt5)/4\approx0.19$ attained at $V=\operatorname{span}\{|00\rangle,(|01\rangle+|12\rangle)/\sqrt2\}$, where
$m_2(V)=(3-\sqrt5)/8$; our argument does not give this. The proof uses two lemmas.

\begin{lemma}\label{lem:qubit}
Let $\mathcal A\ge0$ be an operator on $\mathbb C^2\otimes K$ with blocks $\mathcal A_{jk}$ on $K$, and suppose that
$\Tr_K\mathcal A=I_2$, that is, $\Tr\mathcal A_{00}=\Tr\mathcal A_{11}=1$ and $\Tr\mathcal A_{01}=0$.
\begin{itemize}
\item[(a)] $\lambda_{\max}(\mathcal A)\le1+\tfrac1{\sqrt2}\norm{\mathcal A_{00}-\mathcal A_{11}}_2$.
\item[(b)] Let $\psi$ be a unit vector with $\langle\psi|\mathcal A|\psi\rangle=1+\eta$ and $\eta\ge0$, let $E$ be the
projector onto the support of the reduced operator of $\psi\psi^*$ on $K$, and put
$\varepsilon=\Tr\big((I\otimes(I-E))\mathcal A\big)$. Then
$\norm{\mathcal A_{00}-\mathcal A_{11}}_2^2\ge2\eta(\eta+\varepsilon)$.
\end{itemize}
In particular, if $\lambda=\lambda_{\max}(\mathcal A)\ge1$ and $\tau$ is the sum of all but the two largest eigenvalues
of $\mathcal A_{00}+\mathcal A_{11}$, then $\norm{\mathcal A_{00}-\mathcal A_{11}}_2^2\ge2(\lambda-1)(\lambda-1+\tau)$.
\end{lemma}

The last statement follows from (b) with $\psi$ a top eigenvector, because $E$ has rank at most two, so that
$\varepsilon=\Tr\big((I-E)(\mathcal A_{00}+\mathcal A_{11})\big)\ge\tau$. Part (a) is sharp: equality holds for
$\mathcal A=\lambda\,\varphi\varphi^*+(2-\lambda)\,\varphi'\varphi'^*$ with
$1\le\lambda\le2$, $\varphi=(|0\rangle a+|1\rangle b)/\sqrt2$, $\varphi'=(|0\rangle b-|1\rangle a)/\sqrt2$ and $a\perp b$
unit vectors.

\begin{proof}
Let $\psi$ be a unit vector, with Schmidt decomposition
$\psi=\sqrt p\,q_1\otimes a+\sqrt{1-p}\,q_2\otimes b$, where $a,b\in K$ are orthonormal (enlarging $K$ if
necessary). Conjugating $\mathcal A$ by a diagonal unitary $U$ on $\mathbb C^2$, and replacing $\psi$ by
$(U\otimes I)\psi$, changes neither the hypotheses, nor the spectrum, nor $\mathcal A_{00}-\mathcal A_{11}$, nor
$\langle\psi|\mathcal A|\psi\rangle$, $E$ and $\varepsilon$ in (b). So we may assume that
$q_1=(c,s)$ and, after a phase change of $b$, $q_2=(-s,c)$, where $c=\cos\frac\theta2$, $s=\sin\frac\theta2$ and
$0\le\theta\le\pi$. Put $r=\sqrt{p(1-p)}$ and
\[
Z=\tfrac12I+(p-\tfrac12)\sin\theta\,\sigma_x,\qquad
Y=\tfrac1{\sqrt2}\big[\cos\theta\,(aa^*-bb^*)-2r\sin\theta\,(ab^*+ba^*)\big].
\]
Then $\Tr Z=1$, $Z\ge0$ and $\norm{Y}_2^2=\cos^2\theta+4r^2\sin^2\theta\le1$. We claim that
$\psi\psi^*\le Z\otimes I+\frac1{\sqrt2}\sigma_z\otimes Y$. Both sides commute with $I\otimes E$, where $E$ is the
projector onto $\operatorname{span}\{a,b\}$, so we may check the two blocks separately. On
$\mathbb C^2\otimes\{a,b\}^\perp$ the left side vanishes and the right side is $Z\otimes I\ge0$. On $\mathbb C^2\otimes\operatorname{span}\{a,b\}$, in the basis
$|0a\rangle,|0b\rangle,|1a\rangle,|1b\rangle$, the right side is
\[
P=\begin{pmatrix}c^2&-2rcs&\kappa&0\\-2rcs&s^2&0&\kappa\\\kappa&0&s^2&2rcs\\0&\kappa&2rcs&c^2\end{pmatrix},
\qquad\kappa=(2p-1)cs,
\]
and $\psi=(\sqrt p\,c,\,-\sqrt{1-p}\,s,\,\sqrt p\,s,\,\sqrt{1-p}\,c)$. Using $c^2+s^2=1$ and
$\kappa^2+4r^2c^2s^2=c^2s^2$, one checks that $P^2=P$ and $P\psi=\psi$. So $P$ is a projector whose range contains
$\psi$, and $P-\psi\psi^*\ge0$. Hence
\[
\langle\psi|\mathcal A|\psi\rangle\le\Tr\big(Z\Tr_K\mathcal A\big)
+\tfrac1{\sqrt2}\Tr\big(Y(\mathcal A_{00}-\mathcal A_{11})\big)\le1+\tfrac1{\sqrt2}\norm{\mathcal A_{00}-\mathcal A_{11}}_2 ,
\]
which is (a) when $\psi$ is an eigenvector for $\lambda_{\max}(\mathcal A)$.

For (b) we may assume $\eta>0$. Then $0<p<1$, since $\langle\psi|\mathcal A|\psi\rangle\le1$ when $\psi$ is a product
vector, and $E$ is the projector onto $\operatorname{span}\{a,b\}$. Put $u=|p-\frac12|\sin\theta$ and $\ell=\frac12-u$,
the smallest eigenvalue of $Z$; then $u<\frac12$, so $\ell>0$. Since $4r^2=1-(2p-1)^2$, we have
$\norm{Y}_2^2=\cos^2\theta+4r^2\sin^2\theta=1-4u^2=4\ell(1-\ell)$. Let $\gamma=2\eta/\ell$. By the two blocks above,
$Z\otimes I+\frac1{\sqrt2}\sigma_z\otimes Y-\psi\psi^*$ equals $P-\psi\psi^*\ge0$ on
$\mathbb C^2\otimes\operatorname{span}\{a,b\}$ and $Z\otimes I\ge\ell I$ on $\mathbb C^2\otimes\{a,b\}^\perp$. Hence
$\gamma\big(Z\otimes I+\frac1{\sqrt2}\sigma_z\otimes Y-\psi\psi^*\big)-2\eta\,I\otimes(I-E)\ge0$. Pairing this with
$\mathcal A$ and using $\Tr(Z\Tr_K\mathcal A)=1$ gives
$\frac1{\sqrt2}\Tr\big(Y(\mathcal A_{00}-\mathcal A_{11})\big)\ge\eta+2\eta\varepsilon/\gamma=\eta+\ell\varepsilon$. By the
Cauchy--Schwarz inequality,
\[
\norm{\mathcal A_{00}-\mathcal A_{11}}_2^2\ge\frac{2(\eta+\ell\varepsilon)^2}{\norm{Y}_2^2}
=\frac{(\eta+\ell\varepsilon)^2}{2\ell(1-\ell)}\ge2\eta(\eta+\varepsilon),
\]
where the last step is $(\eta+\ell\varepsilon)^2-4\ell(1-\ell)\eta(\eta+\varepsilon)=\big((1-2\ell)\eta-\ell\varepsilon\big)^2$.
\end{proof}

\begin{lemma}\label{lem:pair}
Let $w_1,w_2$ be an orthonormal basis of a two-dimensional subspace $V\subseteq H_1\otimes H_2$, and let $\Pi$ be as in
the identity above. Then
\[
4\norm{\Pi(w_1\otimes w_2)}^2\ge\tfrac12\min\{\tau_1(V),\tau_2(V)\}.
\]
\end{lemma}

The constant $\frac12$ is best possible. If $\dim H_1=\dim H_2=d\ge3$ and $\sqrt d\,w_1$, $\sqrt d\,w_2$ are unitary
$d\times d$ matrices that are orthogonal in the Hilbert--Schmidt inner product, then the reduced operators $A_j,B_j$ in
the proof below all equal $I/d$, and both sides equal $1-2/d$.

\begin{proof}
For vectors $x_1,x_2$ let $A_j$ and $B_j$ be the reduced operators of $x_jx_j^*$ on $H_1$ and $H_2$. Expanding
$\Pi=\frac14(I-F_1-F_2+F_1F_2)$ gives
\begin{equation}\label{eq:pi4}
4\norm{\Pi(x_1\otimes x_2)}^2=\norm{x_1}^2\norm{x_2}^2+|\langle x_1,x_2\rangle|^2-\Tr A_1A_2-\Tr B_1B_2 .
\end{equation}
Now take $x_j=w_j$ and write $\Delta=4\norm{\Pi(w_1\otimes w_2)}^2$. Then $\rho_1=A_1+A_2$ and $\rho_2=B_1+B_2$ are the
reductions of $P_V$. Let $\mu_1\ge\mu_2\ge\dots$ be the eigenvalues of $\rho_2$. They sum to two, and
$\tau_2=\sum_{i\ge3}\mu_i$.
\begin{itemize}
\item[(i)] By \eqref{eq:pi4}, $\Delta=1-\Tr A_1A_2-\Tr B_1B_2$. Since $\Tr B_j^2=\Tr A_j^2$, this equals
$1-\frac12\Tr\rho_2^2+\frac12\norm{A_1-A_2}_2^2$.
\item[(ii)] Since $\mu_i\le\mu_2$ for $i\ge3$ and $\mu_1+\mu_2+\tau_2=2$,
\[
2-\Tr\rho_2^2-\tau_2=\mu_1(1-\mu_1)+\mu_2(1-\mu_2)-\sum_{i\ge3}\mu_i^2\ge\mu_1(1-\mu_1)+\mu_2(\mu_1-1)
=(1-\mu_1)(\mu_1-\mu_2).
\]
\item[(iii)] The operator $\mathcal A$ on $\mathbb C^2\otimes H_1$ with blocks
$\mathcal A_{jk}=\Tr_2(w_{j+1}w_{k+1}^*)$, $j,k\in\{0,1\}$, is the reduced operator of the vector
$|0\rangle\otimes w_1+|1\rangle\otimes w_2$. So its nonzero eigenvalues are those of $\rho_2$.
Moreover $\Tr_{H_1}\mathcal A=I_2$, $\mathcal A_{00}-\mathcal A_{11}=A_1-A_2$ and
$\mathcal A_{00}+\mathcal A_{11}=\rho_1$. By the last statement of Lemma~\ref{lem:qubit},
$\norm{A_1-A_2}_2^2\ge2(\mu_1-1)(\mu_1-1+\tau_1)$ whenever $\mu_1\ge1$.
\end{itemize}
If $\mu_1\le1$, then (i) and (ii) give $\Delta\ge\frac12\tau_2$. If $\mu_1>1$, then (i)--(iii) and
$\mu_1-\mu_2=2(\mu_1-1)+\tau_2$ give
\[
\Delta\ge\tfrac12\tau_2-\tfrac12(\mu_1-1)(\mu_1-\mu_2)+(\mu_1-1)^2+(\mu_1-1)\tau_1
=\tfrac12(2-\mu_1)\,\tau_2+(\mu_1-1)\,\tau_1 .
\]
The right side is at least $\frac12\mu_1\min\{\tau_1,\tau_2\}\ge\frac12\min\{\tau_1,\tau_2\}$.
\end{proof}

\begin{proof}[Proof of Proposition~\ref{prop:thick}]
Since $q_2(C^*)=q_2(C)$, we may assume that $V$ contains the range of $C$. Write $m=\min\{\tau_1(V),\tau_2(V)\}$; note
$m\le2$. If $\operatorname{rank}C\le1$, then $q_2(C)\ge\frac14\norm{C}_2^2$ by \cite[Prop.~8]{SongChen}, which is
enough. So let
$\operatorname{rank}C=2$, so that $V$ is the range of $C$, and take the decomposition $C=x_1y_1^*+x_2y_2^*$ used for
the identity above. Its Gram matrix $G=(g_{ij})$ has eigenvalues $\sigma_1\ge\sigma_2>0$, and $x_1,x_2$ span $V$. Put
$s=\frac12(\sigma_1+\sigma_2)$ and $\delta=\frac12(\sigma_1-\sigma_2)$. Since every term on the right of the identity is
nonnegative, and $2q_2(C)$ is its left side plus $\frac12(\sigma_1-\sigma_2)^2$,
\[
q_2(C)\ge\delta^2+2\norm{\Pi(x_1\otimes x_2)}^2 .
\]
Write $G=Q^*\operatorname{diag}(\sigma_1,\sigma_2)Q$ with $Q$ unitary. Then
$x_j=\sum_k\sqrt{\sigma_k}\,Q_{kj}v_k$ for an orthonormal basis $v_1,v_2$ of $V$. Let $\hat w_j=\sum_kQ_{kj}v_k$, which
is another orthonormal basis of $V$. Since $\Pi(u\otimes v)=\Pi(v\otimes u)$,
\begin{align*}
\Pi(x_1\otimes x_2)-s\,\Pi(\hat w_1\otimes\hat w_2)
&=\delta Q_{11}Q_{12}\Pi(v_1\otimes v_1)-\delta Q_{21}Q_{22}\Pi(v_2\otimes v_2)\\
&\quad+(\sqrt{\sigma_1\sigma_2}-s)(Q_{11}Q_{22}+Q_{21}Q_{12})\Pi(v_1\otimes v_2).
\end{align*}
By \eqref{eq:pi4}, $\norm{\Pi(v_k\otimes v_k)}^2\le\frac12$ and $\norm{\Pi(v_1\otimes v_2)}^2\le\frac14$. Moreover
$|Q_{11}Q_{12}|,|Q_{21}Q_{22}|\le\frac12$, $|Q_{11}Q_{22}+Q_{21}Q_{12}|\le1$ and $0\le s-\sqrt{\sigma_1\sigma_2}\le\delta$.
So the difference has norm at most $\beta\delta$ with $\beta=\frac1{\sqrt2}+\frac12$, and $4\beta^2=3+2\sqrt2$. By
Lemma~\ref{lem:pair}, $\norm{\Pi(\hat w_1\otimes\hat w_2)}\ge\sqrt\nu$ with $\nu=m/8\le\frac14$. Hence, with
$\zeta=\delta/s\in[0,1]$,
\[
q_2(C)\ge\delta^2+2\big(s\sqrt\nu-\beta\delta\big)_+^2=s^2g(\zeta),\qquad g(\zeta)=\zeta^2+2\big(\sqrt\nu-\beta\zeta\big)_+^2,
\]
while $\norm{C}_2^2=\sigma_1^2+\sigma_2^2=2s^2(1+\zeta^2)$. If $\beta\zeta\ge\sqrt\nu$, then
$g(\zeta)/(2(1+\zeta^2))\ge\zeta^2/(2(1+\zeta^2))\ge\nu/(2\beta^2+2\nu)$, since $\zeta^2/(1+\zeta^2)$ increases with
$\zeta$. If $\beta\zeta<\sqrt\nu$, then $g(\zeta)\ge\min_{z\in\mathbb R}\{z^2+2(\sqrt\nu-\beta z)^2\}=2\nu/(1+2\beta^2)$
and $1+\zeta^2<1+\nu/\beta^2$. Since $2\beta^2+2\nu\le(1+2\beta^2)(1+\nu/\beta^2)$, in both cases
\[
\frac{q_2(C)}{\norm{C}_2^2}\ge\frac{\nu}{(1+2\beta^2)(1+\nu/\beta^2)}
\ge\frac{\nu}{(\frac52+\sqrt2)(4-2\sqrt2)}=\frac{\nu}{6-\sqrt2}=\frac{m}{48-8\sqrt2},
\]
where we used $\nu/\beta^2\le1/(4\beta^2)=3-2\sqrt2$.
\end{proof}

\emph{Range and co-range in one plane.} For these operators, in particular for normal ones, the argument gives Tomar's
inequality with the best possible constant. Tomar proved this constant for every operator whose range lies in a plane
spanned by two vectors with common Schmidt bases \cite[Thm.~7.2]{Tomar2}; here the plane is arbitrary, but the
co-range must lie in it too. If the range and co-range of $C$ lie in a plane $V$, then
\[
q_2(C)\ge\big(1-\tfrac1{\sqrt2}\big)\min_l\lambda_3^{(l)}(V)\,\norm{C}_2^2 ,
\]
and $1-1/\sqrt2>\frac14$. First, $q_2(C)=q_2(\operatorname{Re}C)+q_2(\operatorname{Im}C)$ and
$\norm{C}_2^2=\norm{\operatorname{Re}C}_2^2+\norm{\operatorname{Im}C}_2^2$ for every $C$, because partial traces of
Hermitian operators are Hermitian and $\Tr(PQ)$ is real for Hermitian $P,Q$ (Wu and Zou use the same splitting for
three copies \cite[Cor.~6.2]{WuZou}); so it suffices to treat Hermitian $C$ with range in $V$, and rank one is
immediate. For $C=\omega_1v_1v_1^*+\omega_2v_2v_2^*$ with $v_1,v_2$ orthonormal,
$\sigma_j=|\omega_j|$, $\sigma_1\ge\sigma_2>0$ and $\omega_1\bar\omega_2=\sigma_1\sigma_2e^{i\varphi}$, a direct
computation using \eqref{eq:pi4} gives
\[
q_2(C)=\sigma_1^2\big(\tfrac14+E_1\big)+\sigma_2^2\big(\tfrac14+E_2\big)+\sigma_1\sigma_2\cos\varphi\,\big(\Delta-\tfrac12\big),
\]
where $E_j=1-\Tr\rho_{v_j}^2$ for the reduced operator $\rho_{v_j}$ of $v_jv_j^*$ and
$\Delta=4\norm{\Pi(v_1\otimes v_2)}^2\in[0,1]$. Since $|\cos\varphi\,(\Delta-\frac12)|\le\frac12$, for any $\kappa$
and $\mu$ we get $q_2(C)-\kappa\mu\norm{C}_2^2\ge k_1\sigma_1^2+k_2\sigma_2^2-\frac12\sigma_1\sigma_2$ with
$k_j=\frac14+E_j-\kappa\mu$, which is nonnegative once $k_1,k_2>0$ and $k_1k_2\ge\frac1{16}$ (the two-copy note shows $k_j\ge\frac14-\frac\kappa2>0$). Put
$e_j=\min\{\frac12,E_j\}$. Weyl's inequality $\lambda_3(P+Q)\le\lambda_2(P)+\lambda_2(Q)$ gives
$\mu:=\min_l\lambda_3^{(l)}(V)\le s:=e_1+e_2$, and for $\kappa=1-1/\sqrt2$ the identity
$(\frac14+e_1-\kappa s)(\frac14+e_2-\kappa s)-\frac1{16}=\frac{\sqrt2-1}4s(1-2s)+e_1e_2$ shows $k_1k_2\ge\frac1{16}$
(for $s>\frac12$ use $e_1e_2\ge\frac12(s-\frac12)$). Tomar's example $C=|00\rangle\langle00|-(\sqrt2-1)v_2v_2^*$ with
$v_2=(|11\rangle+|22\rangle)/\sqrt2$ attains the constant. In the same way Ky Fan's maximum principle gives
$\max_l\tau_l(V)\le(1-p_1)+(1-p_2)\le E_1+E_2=:S$ for the largest eigenvalues $p_j$ of $\rho_{v_j}$, and for
$\kappa=(3-\sqrt5)/4$, where $\frac14-\frac\kappa2=\kappa(1-\kappa)$, the identity
$(\frac14+E_1-\kappa S)(\frac14+E_2-\kappa S)-\frac1{16}=\kappa(1-\kappa)S(1-S)+E_1E_2$ (with $E_1E_2\ge S-1$ when
$S>1$) gives $q_2(C)\ge\frac{3-\sqrt5}4\max_l\tau_l(V)\norm{C}_2^2$ on the same class; this is the constant we found
numerically for general $C$. For $C\ge0$, using $\Delta\ge\frac12\min_l\tau_l(V)$ from Lemma~\ref{lem:pair}, a similar
computation gives $q_2(C)\ge\frac14\min_l\tau_l(V)\norm{C}_2^2$, with the tails in place of the third eigenvalues (the two-copy note, Theorem~5(c), gives the full proof). For
indefinite $C$ the third eigenvalues in the first bound cannot be replaced by the tails: for $v_1=|00\rangle$,
$v_2=\sum_{i\ge1}|ii\rangle/\sqrt{d-1}$ and $C=v_1v_1^*-tv_2v_2^*$ the smallest ratio $q_2(C)/(\min_l\tau_l\norm{C}_2^2)$
is $\frac14$ at $d=4$ and $(5-\sqrt{13})/6\approx0.232$ at $d=5$, and it tends to $(3-\sqrt5)/4$ as $d\to\infty$.

\emph{A graded form of the multiplicativity conjecture.} For $0\le k\le n$ put
\[
A_k(C)=\sum_{m=0}^{k}\Big(-\tfrac12\Big)^m\tbinom{n-m}{k-m}S_m(C),\qquad S_0(C)=\norm{C}_2^2 .
\]
Equivalently, $A_k(C)=\sum_{|S|=k}q^{(S)}(C)$, where $q^{(S)}$ is the form \eqref{eq:qn} for the copies in $S$, the
other copies being left untraced. With $u=1-2t$ one checks the identity
\[
\sum_J(-t)^{|J|}\norm{\Tr_JC}_2^2-(1-2t)(1-t)^{n-1}\norm{C}_2^2
=\sum_{k=0}^nu^{n-k}(1-u)^k\Big[A_k(C)-\tbinom{n-1}{k}2^{-k}\norm{C}_2^2\Big].
\]
This is the expansion that Tomar uses to show that the endpoint inequality propagates to smaller $t$
\cite[Lemma~2.3(b)]{Tomar}; there every bracket is replaced by $A_k(C)\ge0$, which gives the weaker bound
$(1-2t)^n\norm{C}_2^2$. The Bernstein weights $u^{n-k}(1-u)^k$ are nonnegative for $0\le t\le1/2$. So the conjectured
identity of the introduction holds on the whole interval as soon as
\[
A_k(C)\ge\tbinom{n-1}{k}2^{-k}\norm{C}_2^2\qquad(k=1,\dots,n,\ \operatorname{rank}C\le2).
\]
For $k=1$ this is the case $m=1$ of \eqref{eq:CS}, proved here, and for $k=n$ it is \eqref{eq:qn}. For $n=2$ both
levels are known, the second being the two-copy theorem, so the identity holds for two copies at every
$t\in[0,1/2]$. For $n=3$ the middle level, $11S_0-4S_1+S_2\ge0$, also follows from known results. One checks
\[
11S_0-4S_1+S_2=\tfrac56\cdot8q_3+\tfrac23\sum_{i=1}^3\Big[2S_0+\tfrac12S_3-\norm{\Tr_iC}_2^2-\norm{\Tr_{jk}C}_2^2\Big]
+\tfrac16\big(2S_0-S_3\big),
\]
where $\{j,k\}$ is the complement of $i$ and $S_3=|\Tr C|^2$. The bracket is the slack of the two-copy theorem
\cite{FHPV,BGH} for the bipartition of the three factors into $\{i\}$ and $\{j,k\}$, and $|\Tr C|^2\le2\norm{C}_2^2$
for rank two. So, given the three-copy theorem \cite{Tomar}, all three levels hold for $n=3$, and the conjectured
identity of the introduction holds for three copies at every $t\in[0,1/2]$. For local dimension at most four this three-copy statement is due to MacCormack and Bao \cite[Thm.~1.2]{MB}, whose Proposition~3.1 is an exact identity of the same kind; what is new here is every $d$, given Tomar's theorem, the graded formulation for every $n$, and the linear program below. For $n\ge4$ no such linear combination
exists: a linear program over the two- and three-copy theorems applied to all groupings of the copies (with the other
copies as spectators), the bounds \eqref{eq:CS}, Theorem~\ref{thm:main} and the bound with a flatness term above finds
none for any level $2\le k\le n$ at $n=4,5$. In our numerics every level holds, with equality at
$C=P\otimes|c\rangle\langle c|^{\otimes(n-1)}$ for a rank-two projector $P$ on one factor, for all $n\le6$ at $d=3$,
with spot checks at $d=2$ and $d=4$. Summing the two-copy inequality over all pairs of copies gives only
$A_2(C)\ge0$, which misses the required constant.

\emph{Below the PPT point.} For $0\le t\le1/d$ the form $\sum_J(-t)^{|J|}\norm{\Tr_JZ}_2^2$ is positive semidefinite in
$Z$, and there a rank-free version of the conjectured identity seems to hold: for every operator $Z$ on
$(\mathbb C^d)^{\otimes n}$,
\[
\sum_J(-t)^{|J|}\norm{\Tr_JZ}_2^2\ge(1-t)^{n-1}\big(\norm{Z}_2^2-t\norm{Z}_1^2\big)\qquad(0\le t\le1/d).
\]
Its derivative at $t=0$ is Theorem~\ref{thm:main}, and it is an equality for $P\otimes|c\rangle\langle c|^{\otimes(n-1)}$
with $P\ge0$. For rank two and $d=3$ it would give the conjectured identity of the introduction on $[0,1/3]$, for every
$n$. At $t=1/d$ it says that the part of $Z$ that is traceless on every factor has squared norm at least
$(1-1/d)^{n-1}(\norm{Z}_2^2-\norm{Z}_1^2/d)$. In our searches it holds for $(d,n)=(2,2),(2,3),(2,4)$ at $t=1/2$, for
$(3,2)$ and $(3,3)$ at several $t\le1/3$, and for $(4,2)$ at $t=1/5,1/4$. It fails as soon as $t>1/d$: for $d=3$ and
$n=2$ already at $t=0.34$, with $Z$ of rank three. The one-site induction behind Theorem~\ref{thm:main} reduces it to a
block inequality, which holds at $t=0$ (it is Lemma~\ref{lem:R}) but fails at $t=1/d$ in a fixed block basis.

\emph{Other Tsallis indices.} The bound $\sum_i\Tr[(\Tr_i\rho)^q]\le(n-1)\Tr\rho^q+(\Tr\rho)^q$ also holds in our
numerics for $q=3/2$ and $q=3$, at every rank. A Han inequality for all Tsallis entropies with $q>1$ would therefore be
natural. For $n=2$ it is Audenaert's theorem \cite{Aud}. Our argument is specific to $q=2$.

\section{What the method does not give}
As in Section~\ref{sec:remarks}, the numerical statements in this section come from exploratory searches that are not part of the scripted checks.


The endpoint inequality \eqref{eq:qn} needs upper bounds on the odd moments $S_m$ and lower bounds on the even ones.
Corollary~\ref{cor:allm} supplies only the first kind. This is why \cite{CS} needed a further unproved assumption to
reach undistillability.

The trace-norm mechanism behind Theorem~\ref{thm:main} also does not extend directly. We tested three natural rank-free
forms of \eqref{eq:qn} numerically. All of them fail already at $d=3$:
\begin{itemize}
\item $\Tr[(Z\otimes Z^*)(2(-1)^N-(-3)^N)]\le\norm{Z}_1^2$, where $N$ counts the locally antisymmetric factors of two
copies. It fails for $n=2$.
\item $q_n(X)\ge0$ for all $X$ with $\norm{X}_1^2\le2\norm{X}_2^2$. It fails for $n=2$.
\item $q_n(\rho)\ge0$ for all density matrices with $\Tr\rho^2\ge1/2$. It fails for $n=3$.
\end{itemize}
So the rank-two constraint has to be used in an essential way at the endpoint. At first order it does not: the bound of
Corollary~\ref{cor:rank} only sees the rank through $\norm{C}_1^2\le r\norm{C}_2^2$.

The same contrast appears for positive maps. Inequality \eqref{eq:qn} for all rank-two $C$ says that
$\mathrm{id}_2\otimes\psi^{\otimes n}$ is positive, where $\psi(X)=\Tr(X)I-X/2$. Corollary~\ref{cor:map} says that the
first-order part of this map stays positive with an ancilla of any dimension.

\begin{thebibliography}{99}
\bibitem{DSSTT} D. P. DiVincenzo, P. W. Shor, J. A. Smolin, B. M. Terhal, A. V. Thapliyal, Phys. Rev. A \textbf{61}, 062312 (2000).
\bibitem{DCLB} W. D\"ur, J. I. Cirac, M. Lewenstein, D. Bru\ss, Phys. Rev. A \textbf{61}, 062313 (2000).
\bibitem{HRZ} P. Horodecki, \L. Rudnicki, K. \.Zyczkowski, Five open problems in quantum information theory, PRX Quantum \textbf{3}, 010101 (2022).
\bibitem{FGP} J. Fu, L. Gao, S.-J. Park, arXiv:2607.21367.
\bibitem{SongChen} Z. Song, L. Chen, arXiv:2607.23416.
\bibitem{FHPV} T. C. Fraser, F. Huber, B. Pozsgay, I. Vona, On the two-copy distillability of Werner states and a new partial trace inequality, arXiv:2607.24309.
\bibitem{BGH} K. Bharti, R. Gajjala, T. Haug, Two-copy nondistillability of Werner states: sharp partial-trace inequalities and finite-copy extensions, arXiv:2607.24479.
\bibitem{WuZou} T. Wu, Q. Zou, Sharp Pl\"ucker geometry for three-copy Werner distillation, arXiv:2608.02647.
\bibitem{Tomar} T. S. Tomar, Undistillability of Werner states from three, four and five copies, Zenodo 23718035 (2026).
\bibitem{Tomar2} T. S. Tomar, Equality cases and every-copy results for the Werner endpoint inequality, Zenodo 23718037 (2026).
\bibitem{CR23} P. Costa Rico, New partial trace inequalities and distillability of Werner states, Lett. Math. Phys. \textbf{115}, 47 (2025); arXiv:2310.05726.
\bibitem{QGWG} S.-Y. Qi, G. Gupur, Y.-C. Wu, G.-P. Guo, Unveiling NPT bound problem: from distillability sets to inequalities and multivariable insights, Phys. Rev. A \textbf{110}, 012406 (2024); arXiv:2402.18037.
\bibitem{AnonD3} Anonymous, Three-copy undistillability of the boundary qutrit Werner state (2026), Zenodo 10.5281/zenodo.23134454.
\bibitem{LiuD4} A.-X. Liu, C.-F. Qiao, Three-copy undistillability of the boundary two-ququart Werner state (2026), github.com/GLiuax/d4-three-copy-werner.
\bibitem{MB} I. MacCormack, N. Bao, Three copies of the Werner state: the sharp curve in dimensions three and four and a bound in every dimension (2026), github.com/physical-superintelligence/oqp-batch-2026.
\bibitem{CS} I. Chattopadhyay, D. Sarkar, NPT bound entanglement---the problem revisited, arXiv:quant-ph/0609050.
\bibitem{PPHH} \L. Pankowski, M. Piani, M. Horodecki, P. Horodecki, A few steps more towards NPT bound entanglement, IEEE Trans. Inf. Theory \textbf{56}, 4085 (2010); arXiv:0711.2613.
\bibitem{Aud} K. M. R. Audenaert, Subadditivity of $q$-entropies for $q>1$, J. Math. Phys. \textbf{48}, 083507 (2007).
\bibitem{Ras} A. E. Rastegin, Relations for certain symmetric norms and anti-norms before and after partial trace, J. Stat. Phys. \textbf{148}, 1040 (2012).
\bibitem{CRW} P. Costa Rico, M. M. Wolf, Partial trace relations beyond normal matrices, J. Math. Phys. \textbf{67}, 082203 (2026); arXiv:2507.18278.
\bibitem{FYK} S. Furuichi, N. Minculete, F.-C. Mitroi, Some inequalities on generalized entropies, J. Inequal. Appl. \textbf{2012}, 226 (2012); arXiv:1104.0360.
\end{thebibliography}

\end{document}
